\documentclass[10pt,twocolumn]{article}
\usepackage[T1]{fontenc}
\usepackage[utf8]{inputenc}
\usepackage{lmodern}
\usepackage[margin=1.8cm,top=2.4cm,bottom=2cm,columnsep=0.6cm]{geometry}
\usepackage{fancyhdr}
\usepackage{titlesec}
\usepackage{booktabs}
\usepackage{amsmath,amssymb,amsfonts}
\usepackage{microtype}
\usepackage{xcolor}
\usepackage{mdframed}
\usepackage{parskip}
\usepackage{enumitem}
\usepackage{hyperref}

\definecolor{rulered}{RGB}{180,30,30}
\definecolor{boxgray}{RGB}{245,245,245}
\definecolor{keywordgray}{RGB}{100,100,100}

\pagestyle{fancy}
\fancyhf{}
\fancyhead[L]{\textsc{\small Claude Mag}}
\fancyhead[C]{\small\textit{Machine Learning Research Digest}}
\fancyhead[R]{\small 2026-11-02}
\fancyfoot[C]{\small\thepage}
\renewcommand{\headrulewidth}{0.8pt}

\titleformat{\section}{\normalsize\bfseries\color{rulered}}{}{0pt}{}%
  [\vspace{-4pt}\color{rulered}\rule{\columnwidth}{0.4pt}]
\titlespacing{\section}{0pt}{8pt}{4pt}

\hypersetup{colorlinks=true,urlcolor=rulered,linkcolor=black}

\begin{document}
\setlength{\parindent}{0pt}
\setlength{\parskip}{4pt}

{\small\color{keywordgray}\textbf{Keywords:}
  multi-agent LLMs \textbullet{}
  collective reasoning \textbullet{}
  alignment measurement \textbullet{}
  information-efficient communication}\par\vspace{4pt}

{\LARGE\bfseries\raggedright
  Collective Emergent Alignment in Multi-Agent LLMs}\par\vspace{4pt}

{\large\itshape\raggedright
  REMAP Quantifies Interaction-Induced Reasoning Gains;
  Gap-Aware Factual Messages Amplify Them by $+3.1$ F1
  While Cutting Token Cost by 58\%}\par\vspace{6pt}

{\small\textbf{By} Boxuan Wang, Tianxiang Sun, Linyang Li,
  Xipeng Qiu, Hang Yan}\par\vspace{2pt}

{\small\color{keywordgray}arXiv:2610.34836 \textbullet{} cs.AI
  \textbullet{} submitted 26 Oct 2026}
\par\vspace{2pt}\noindent{\color{rulered}\rule{\columnwidth}{0.6pt}}\vspace{6pt}

Multi-agent LLM systems often outperform any individual agent on
reasoning tasks---but is this due to genuine peer interaction or simply
to extra self-refinement compute? This paper introduces Collective
Emergent Alignment (CEA) to measure the interaction-specific gain, and
the REMAP framework to quantify it at the level of individual facts.
Across seven model families and two multi-hop benchmarks, CEA is
consistently $+6.3$ F1 points above matched self-refinement, and
gap-aware factual communication amplifies this by a further $+3.1$
points while reducing communication tokens by 58\%.

\begin{mdframed}[backgroundcolor=boxgray,linecolor=rulered,linewidth=1pt,
    innerleftmargin=6pt,innerrightmargin=6pt,
    innertopmargin=6pt,innerbottommargin=6pt]
\textbf{\small AT A GLANCE}\par\vspace{4pt}
\begin{description}[leftmargin=0pt,labelsep=4pt,itemsep=2pt,parsep=0pt,
    font=\small\bfseries]
  \item[Problem:] It is unclear whether multi-agent LLM performance
    gains reflect genuine collective intelligence or just additional
    self-refinement compute.
  \item[Why it matters:] Multi-agent LLM systems are increasingly
    deployed in production; understanding whether interaction
    actually helps is essential for designing efficient systems.
  \item[Method:] REMAP reference-entailment matrix measures per-fact
    alignment gains attributable to interaction; gap-aware factual
    messages route only missing facts between agents.
  \item[Result:] CEA is $+6.3$ F1 over matched self-refinement;
    gap-aware communication adds $+3.1$ F1 and cuts tokens by 58\%.
  \item[Evidence:] 7 model families $\times$ 2 multi-hop QA benchmarks
    (MuSiQue, 2WikiMultiHopQA); group sizes 2, 4, 8.
\end{description}
\end{mdframed}

\section{The Big Picture}

In a multi-agent LLM system, several models communicate over multiple
rounds to jointly arrive at an answer. The intuition is that diverse
agents will catch each other's errors and combine partial knowledge.
Empirically, these systems often outperform single-model baselines.
But this comparison is confounded: a multi-agent system at round $k$
has performed more computation than a single agent that runs $k$
self-refinement steps on its own. The improvement could be
\emph{collective}---arising from information transfer between agents---
or \emph{individual}---arising from additional self-refinement
at matched compute.

This paper disentangles the two. It defines Collective Emergent
Alignment (CEA) as the improvement in reference-grounded reasoning
that arises specifically from inter-agent interaction, over and above
what the same total compute would yield in independent self-refinement.

\section{Why Existing Approaches Fall Short}

\textbf{Multi-agent vs.\ single-agent comparisons:} Most prior work
compares a multi-agent system directly to a single-agent baseline
without controlling for compute. The multi-agent system always
has more iterations and more tokens; any gain is ambiguous.

\textbf{Aggregate metrics only:} Existing evaluations report final
F1 or accuracy but do not attribute performance gains to specific
communication events. There is no way to tell which round or which
message was responsible for a correct answer.

\textbf{No formal model of collective gain:} Prior work lacks a
framework for defining what ``collective'' means in a measurable
sense. Terms like ``debate,'' ``collaboration,'' and ``consensus''
are used loosely without formal grounding.

\section{Core Mechanism}

\textbf{Reference-Entailment Matrix.} REMAP constructs a binary
matrix $M \in \{0, 1\}^{A \times R \times F}$ where $A$ is the
number of agents, $R$ the number of rounds, and $F$ the number of
atomic facts in the reference answer. $M_{a, r, f} = 1$ if fact $f$
is entailed by agent $a$'s response at round $r$. The matrix is
populated using a pretrained natural-language inference model.

\textbf{Collective Emergent Alignment.} A fact $f$ is attributed to
CEA if:
\begin{itemize}[itemsep=1pt,parsep=0pt]
  \item It first appears in some agent's response after a communication
    round (not at round 0), and
  \item It does not appear in the counterfactual trajectory where
    that agent self-refines without receiving peer messages.
\end{itemize}
The CEA score is the fraction of reference facts attributed to CEA
under this definition.

\textbf{Gap-Aware Factual Communication.} Instead of exchanging full
reasoning traces, each agent identifies which facts in its peer's
last response are absent from or contradict its own response.
It then routes a compact message containing only those ``gap facts''
to the peer. This is implemented with a lightweight LLM step that
extracts the gap-fact set from both responses.

\section{Technical Formulation}

Let $R_{a, r}$ denote agent $a$'s response at round $r$, and let
$E(R, f)$ be the entailment probability for fact $f$ given response
$R$. The binary matrix is:
\[
M_{a,r,f} = \mathbf{1}[E(R_{a,r}, f) > \tau]
\]
The Semantic Coverage of the collective at round $r$ is:
\[
\text{SC}(r) = \frac{|\{f : \exists\, a,\; M_{a,r,f} = 1\}|}{|F|}
\]
The CEA score is:
\[
\text{CEA} = \frac{1}{|F|}\sum_f \mathbf{1}[f \in F_{\text{collective}}
  \setminus F_{\text{individual}}]
\]
where $F_{\text{collective}}$ are facts recovered by the multi-agent
system and $F_{\text{individual}}$ are those recovered by the
matched self-refinement baseline.

For gap-aware communication, the message from agent $a$ to agent $b$
at round $r$ is:
\[
\mathcal{G}_{a \to b}^r = \{f : M_{a,r,f} = 1,\; M_{b,r,f} = 0\}
\]
and the message content is a natural-language summary of
$\mathcal{G}_{a \to b}^r$ generated by a compact auxiliary LLM.

\section{Learning or Inference Procedure}

No model training is required for REMAP measurement. All LLMs are
used in their pretrained, instruction-tuned states. The evaluation
procedure is:
\begin{enumerate}[itemsep=2pt,parsep=0pt]
  \item Initialize $A$ agents with the same question.
  \item Run $R$ rounds of communication, alternating between
    standard full-trace and gap-aware factual modes.
  \item Construct $M$ using the NLI model after each round.
  \item Compute SC, CEA, and per-round alignment gains.
\end{enumerate}

The self-refinement counterfactual is obtained by re-running each
agent independently for the same number of rounds, receiving no
peer messages.

\section{What the Guarantee Says}

The paper provides no formal guarantee about CEA's optimality or
about the gap-aware communication protocol. The theoretical
contribution is the definition of CEA itself, which is grounded
in the entailment matrix and admits a clean counterfactual
interpretation. The gap-aware communication protocol is justified
empirically: it consistently amplifies CEA and reduces token cost
across all evaluated settings.

\section{Experimental Findings}

\begin{itemize}[itemsep=2pt,parsep=0pt]
  \item \textbf{CEA is real:} Across 7 model families and both
    benchmarks, CEA averages \textbf{$+6.3$ F1 points}
    over matched self-refinement. All model families show
    positive CEA; none is negative.
  \item \textbf{Gap-aware communication:} Amplifies CEA by
    \textbf{$+3.1$ F1} while reducing communication tokens
    by \textbf{58\%} compared to full-trace exchange.
  \item \textbf{Group size:} CEA increases with group size (2 to 8
    agents) but with diminishing returns; going from 2 to 4 agents
    yields $+4.1$ F1, while going from 4 to 8 yields only $+1.8$ F1.
  \item \textbf{Benchmarks:} MuSiQue (harder, requiring 4-hop
    reasoning) shows higher CEA ($+8.1$) than 2WikiMultiHopQA
    ($+4.5$), consistent with more room for collective benefit when
    individual models struggle.
\end{itemize}

\section{Ablations and Interpretation}

\textbf{Communication timing:} Communicating every round is better
than every 2 rounds ($+1.2$ F1), but the gap narrows with gap-aware
messages (which are shorter and can be communicated more frequently).

\textbf{Message content:} Full-trace exchange vs.\ gap-aware facts:
full-trace exchange achieves 3.1\% lower F1 and costs 58\% more
tokens. Gap-aware messages are both cheaper and better.

\textbf{NLI threshold $\tau$:} CEA is robust to $\tau \in [0.6, 0.8]$;
below 0.6, noisy entailment decisions inflate CEA artificially;
above 0.8, the definition becomes too strict and CEA underestimates
the true gain.

\textbf{Dominance phenomenon:} In 22\% of episodes, one agent
provides $>$70\% of the CEA facts. Sequential communication
reduces this to 9\%, suggesting that synchronous communication
can lead to one agent dominating the group.

\section*{Reference}

Boxuan Wang, Tianxiang Sun, Linyang Li, Xipeng Qiu, Hang Yan.
\textbf{Collective Emergent Alignment: Measuring and Amplifying
Interaction-Induced Reasoning Alignment in Multi-Agent LLMs.}
arXiv:2610.34836. 2026.
\url{https://arxiv.org/abs/2610.34836}

\end{document}
